\documentclass[11pt]{article}
\usepackage{mathblatt}
% gesetzt von werkzeuge/setzer.py aus dem Lernweg (L3-1 · L3-2 · L3-3 · L3-4 · L3-5 · L3-6 · L3-7) und den Bankzeilen; Muster T6B.
% Präambel des Setzers (werkzeuge/setzer.py), wortgleich aus dem Hand-Satz T6B (bau/proben/2026-10-09/kathete-berechnen), dort aus K4W.
\makeatletter
\setlength{\mbpfbe}{0mm}% keine Punkte (Bauregel 6.4)
% eingerückt (Bauregel 4.7, 6.6): \gEin Einzug, \gSchrift eine Stufe kleiner
\newlength{\gEin}\setlength{\gEin}{0pt}
\let\gSchrift\relax
\renewcommand{\pfkreuzzeile}[1]{\par\noindent\hangindent3.2em\hangafter1\hspace*{1.6em}$\square$\ #1\par}
\newcommand{\gkreuz}[1]{$\square$\ #1\hspace{2em}}
% Bauregel 6.7: links, was man ansieht, rechts, was man tut; Spaltengrenze überall an derselben Stelle
\newsavebox{\gbox}\newlength{\gLinks}
\newcommand{\gzwei}[2]{\par\smallskip\noindent\sbox{\gbox}{#1}%
  \setlength{\gLinks}{\dimexpr0.5\textwidth-\gEin-\mbpfnr\relax}%
  \ifdim\wd\gbox>\dimexpr\gLinks-4mm\relax
    \usebox{\gbox}\par\smallskip\noindent\begin{minipage}{\linewidth}#2\end{minipage}%
  \else
    \begin{minipage}[t]{\gLinks}\vspace{0pt}\usebox{\gbox}\end{minipage}%
    \begin{minipage}[t]{\dimexpr\linewidth-\gLinks\relax}\vspace{0pt}\raggedright #2\end{minipage}%
  \fi\par}
\RenewDocumentEnvironment{pfaufg}{m m m}{%
  \begin{lrbox}{\mbaufgabebox}\begin{minipage}[t]{\linewidth}%
  \gSchrift\noindent\hspace*{\gEin}\mbpfrandmarke{#1}%
  \begin{minipage}[t]{\mbpfnr}\strut\textbf{#2}\end{minipage}%
  \begin{minipage}[t]{\dimexpr\linewidth-\gEin-\mbpfnr\relax}\raggedright\strut\ignorespaces}%
 {\end{minipage}%
  \end{minipage}\end{lrbox}%
  \setlength{\mbaufgabehoehe}{\dimexpr\ht\mbaufgabebox+\dp\mbaufgabebox\relax}%
  \par\addvspace{7pt}\Needspace*{\mbaufgabehoehe}%
  \noindent\usebox{\mbaufgabebox}\par\addvspace{7pt}}
\makeatother
% Rechenraum als Karo (Bauregel 6.9): n Rechenschritte = n mal 10 mm
\newcommand{\karo}[1]{\par\smallskip\noindent\begin{tikzpicture}
  \clip (0,0) rectangle ({\linewidth-1mm},{#1*10mm});
  \draw[black!16,line width=0.3pt,step=5mm] (0,0) grid ({\linewidth},{#1*10mm});
  \end{tikzpicture}\par}
\newcommand{\kk}[1]{$\square$\,#1\hspace{0.9em}}
\newcommand{\linie}[1]{\,{\color{black!45}\rule[-1pt]{#1}{0.4pt}}\,}
% Zeile am Ende jedes Durchgangs: Originale klein grau, darüber; dann Vorgänger/Nachfolger (Bauregel 3.4, 6.5)
\newcommand{\blattende}[2]{\par\vfill\noindent\begin{minipage}{\linewidth}{\scriptsize\color{mbgrau}Originale: #1\par}\smallskip
{\footnotesize #2\par}\end{minipage}}
\newcommand{\pkt}{\textperiodcentered{}}

\newcommand{\kennung}{UV3}
\newcommand{\grau}[1]{{\color{mbgrau}#1}}

\begin{document}
\pfheftstil
\fancyfoot[C]{\parbox[t]{\textwidth}{\mbfhbox\par\vspace{1.5mm}{\footnotesize\color{mbgrau}{\scriptsize \kennung}\hfill\thepage}}}
\pfheftkopf{Das Dreieck erst finden}{Klasse 9}
% L3-1: In jeder Figur und jedem Körper steckt ein rechtwinkliges Dreieck; man findet es am rechten Winkel und nennt die Hypotenuse gegenüber
\begin{pfaufg}{}{1.}{}
Jede Figur enthält ein rechtwinkliges Dreieck. Fahre es farbig nach. Welche Strecke ist seine Hypotenuse?\par\medskip\noindent
\begin{minipage}[b]{0.32\linewidth}\centering
\begin{tikzpicture}[line width=0.7pt]\coordinate (A) at (0,0);\coordinate (B) at (3.2,0);\coordinate (C) at (2.4,1.6);\coordinate (D) at (0.8,1.6);\coordinate (F) at (0.8,0);\draw (A)--(B)--(C)--(D)--cycle;\draw[dashed,line width=0.5pt] (D)--(F);\mbrwmarke{(F)}{(A)}{(D)}\node[font=\small,anchor=north] at (A) {$A$};\node[font=\small,anchor=north] at (B) {$B$};\node[font=\small,anchor=south] at (C) {$C$};\node[font=\small,anchor=south] at (D) {$D$};\node[font=\small,anchor=north] at (F) {$F$};\end{tikzpicture}\par\smallskip
{\small a)\ Trapez}\par\smallskip
\linie{30mm}
\end{minipage}\hfill
\begin{minipage}[b]{0.32\linewidth}\centering
\begin{tikzpicture}[line width=0.7pt]\coordinate (M) at (0,0);\coordinate (P) at (1.2,0);\coordinate (S) at (0,2.6);\draw (-1.2,0)--(S)--(P);\draw (-1.2,0) arc (180:360:1.2 and 0.4);\draw[dashed,line width=0.5pt] (-1.2,0) arc (180:0:1.2 and 0.4);\draw[dashed,line width=0.5pt] (M)--(S) (M)--(P);\fill (M) circle (0.8pt);\mbrwmarke{(M)}{(P)}{(S)}\node[font=\small,anchor=south] at (S) {$S$};\node[font=\small,anchor=north east] at (M) {$M$};\node[font=\small,anchor=west] at (P) {$P$};\end{tikzpicture}\par\smallskip
{\small b)\ Kegel}\par\smallskip
\linie{30mm}
\end{minipage}\hfill
\begin{minipage}[b]{0.32\linewidth}\centering
\begin{tikzpicture}[line width=0.7pt]\coordinate (A) at (0,0);\coordinate (B) at (3,0);\coordinate (C) at (4,0.8);\coordinate (D) at (1,0.8);\coordinate (M) at (2,0.4);\coordinate (S) at (2,2.7);\coordinate (H) at (1.5,0);\draw (A)--(B)--(C)--(S)--cycle (S)--(B);\draw[dashed,line width=0.5pt] (A)--(D)--(C) (D)--(S) (S)--(M)--(H);\draw[line width=0.5pt] (S)--(H);\fill (M) circle (0.8pt);\mbrwmarke{(M)}{(H)}{(S)}\node[font=\small,anchor=south] at (S) {$S$};\node[font=\small,anchor=west] at (M) {$M$};\node[font=\small,anchor=north] at (H) {$H$};\end{tikzpicture}\par\smallskip
{\small c)\ Pyramide}\par\smallskip
\linie{30mm}
\end{minipage}
\fusshilfe{\mbox{1:~a) AD; b) SP; c) SH}}
\end{pfaufg}

% L3-2: Ganz oder halb? Die Höhe trifft die Mitte (halbe Länge); ein Stab von Rand zu Rand nimmt die ganze
\begin{pfaufg}{}{2.}{}
Im Kegel und in der Pyramide gehört die Höhe zu einem rechtwinkligen Dreieck, im Becher der Stab. Welche Länge ist dort die waagerechte Kathete? Kreuze an.\par\medskip\noindent
\begin{minipage}[b]{0.32\linewidth}\centering
\begin{tikzpicture}[line width=0.7pt]\coordinate (M) at (0,0);\coordinate (P) at (1.2,0);\coordinate (S) at (0,2.6);\draw (-1.2,0)--(S)--(P);\draw (-1.2,0) arc (180:360:1.2 and 0.4);\draw[dashed,line width=0.5pt] (-1.2,0) arc (180:0:1.2 and 0.4);\draw[dashed,line width=0.5pt] (M)--(S) (M)--(P);\fill (M) circle (0.8pt);\mbrwmarke{(M)}{(P)}{(S)}\draw[|-|,line width=0.4pt] (-1.2,-0.65)--(1.2,-0.65) node[midway,below,font=\small] {$18$ cm};\end{tikzpicture}\par\smallskip
{\small a)\ \kk{$9$ cm}\par\hspace*{1.1em}\kk{$18$ cm}}
\end{minipage}\hfill
\begin{minipage}[b]{0.32\linewidth}\centering
\begin{tikzpicture}[line width=0.7pt]\coordinate (A) at (0,0);\coordinate (B) at (3,0);\coordinate (C) at (4,0.8);\coordinate (D) at (1,0.8);\coordinate (M) at (2,0.4);\coordinate (S) at (2,2.7);\coordinate (H) at (1.5,0);\draw (A)--(B)--(C)--(S)--cycle (S)--(B);\draw[dashed,line width=0.5pt] (A)--(D)--(C) (D)--(S) (S)--(M)--(H);\draw[line width=0.5pt] (S)--(H);\fill (M) circle (0.8pt);\mbrwmarke{(M)}{(H)}{(S)}\node[font=\small,anchor=north] at (1.5,-0.05) {$7$ m};\end{tikzpicture}\par\smallskip
{\small b)\ \kk{$3{,}5$ m}\par\hspace*{1.1em}\kk{$7$ m}}
\end{minipage}\hfill
\begin{minipage}[b]{0.32\linewidth}\centering
\begin{tikzpicture}[line width=0.7pt]\draw (0,2.4) ellipse (0.8 and 0.28);\draw (-0.8,0)--(-0.8,2.4) (0.8,0)--(0.8,2.4);\draw (-0.8,0) arc (180:360:0.8 and 0.28);\draw[dashed,line width=0.5pt] (-0.8,0) arc (180:0:0.8 and 0.28);\draw[line width=1.6pt] (-0.8,0)--(1.1,2.85);\draw[dashed,line width=0.5pt] (0,0)--(0.8,0);\fill (0,0) circle (0.8pt);\node[font=\small,anchor=north] at (0.4,-0.3) {$r = 3$ cm};\end{tikzpicture}\par\smallskip
{\small c)\ \kk{$3$ cm}\par\hspace*{1.1em}\kk{$6$ cm}}
\end{minipage}
\fusshilfe{\mbox{2:~a) 9 cm; b) 3,5 m; c) 6 cm}}
\end{pfaufg}

% L3-3: Im gleichschenkligen Dreieck erst halbieren, dann entscheiden: Höhe gesucht (minus) oder Schenkel gesucht (plus)
\begin{pfaufg}{}{3.}{}
Die Figur zeigt ein gleichschenkliges Dreieck. Fahre zuerst das rechtwinklige Dreieck nach.
\gzwei{\begin{tikzpicture}[line width=0.7pt]\coordinate (A) at (0,0);\coordinate (B) at (3,0);\coordinate (C) at (1.5,2.2);\coordinate (F) at (1.5,0);\draw (A)--(B)--(C)--cycle;\draw[dashed,line width=0.5pt] (C)--(F);\mbrwmarke{(F)}{(B)}{(C)}\node[font=\small,anchor=north] at (1.5,-0.05) {$g$};\node[font=\small,anchor=south west] at (2.3,1.15) {$s$};\node[font=\small,anchor=south east] at (0.7,1.15) {$s$};\node[font=\small,anchor=west] at (1.5,0.9) {$h$};\end{tikzpicture}}{%
\textbf{a)}\ $g = 30$ cm, $s = 17$ cm. Berechne $h$.\linie{14mm}cm\par\bigskip
\textbf{b)}\ $g = 42$ cm, $h = 20$ cm. Berechne $s$.\linie{14mm}cm\karo{5}}
\fusshilfe{\mbox{3:~a) 8 cm; b) 29 cm}}
\end{pfaufg}

% L3-4: Dasselbe im Körper: Kegel mit Radius aus dem Durchmesser, Mantellinie oder Höhe gesucht
\begin{pfaufg}{}{4.}{}
Die Figur zeigt einen Kegel. Fahre zuerst das rechtwinklige Dreieck nach.
\gzwei{\begin{tikzpicture}[line width=0.7pt]\coordinate (M) at (0,0);\coordinate (P) at (1.2,0);\coordinate (S) at (0,2.6);\draw (-1.2,0)--(S)--(P);\draw (-1.2,0) arc (180:360:1.2 and 0.4);\draw[dashed,line width=0.5pt] (-1.2,0) arc (180:0:1.2 and 0.4);\draw[dashed,line width=0.5pt] (M)--(S) (M)--(P);\fill (M) circle (0.8pt);\mbrwmarke{(M)}{(P)}{(S)}\draw[|-|,line width=0.4pt] (-1.2,-0.65)--(1.2,-0.65) node[midway,below,font=\small] {$d$};\node[font=\small,anchor=east] at (0,1.3) {$h$};\node[font=\small,anchor=south west] at (0.7,1.4) {$s$};\end{tikzpicture}}{%
\textbf{a)}\ $d = 10$ cm, $h = 12$ cm. Berechne $s$.\linie{14mm}cm\par\bigskip
\textbf{b)}\ $s = 29$ cm, $d = 40$ cm. Berechne $h$.\linie{14mm}cm\karo{5}}
\fusshilfe{\mbox{4:~a) 13 cm; b) 21 cm}}
\end{pfaufg}

% L3-5: Strecke im Koordinatensystem: das Dreieck selbst einzeichnen, Katheten als Unterschiede abzählen, auch über die Achse
\begin{pfaufg}{}{5.}{}
Im Koordinatensystem sind die Punkte $A$ und $B$ eingezeichnet.
\gzwei{\begin{tikzpicture}[x=0.6cm,y=0.6cm,line width=0.7pt]\draw[black!18,line width=0.3pt] (-4,-3) grid (3,5);\draw[-{Stealth[length=4pt]},line width=0.5pt] (-4,0)--(3.4,0) node[right,font=\small] {$x$};\draw[-{Stealth[length=4pt]},line width=0.5pt] (0,-3)--(0,5.4) node[above,font=\small] {$y$};\foreach \x in {-3,-2,-1,1,2} \node[font=\scriptsize,below] at (\x,0) {$\x$};\foreach \y in {-2,-1,1,2,3,4} \node[font=\scriptsize,left] at (0,\y) {$\y$};\fill (-3,-2) circle (2.2pt) node[below left,font=\small] {$A$};\fill (2,4) circle (2.2pt) node[above left,font=\small] {$B$};\draw[line width=0.9pt] (-3,-2)--(2,4);\end{tikzpicture}}{%
\textbf{a)}\ Zeichne das rechtwinklige Dreieck mit einer waagerechten und einer senkrechten Kathete ein. Wie lang sind die Katheten (in Kästchen)?\linie{14mm}\par\bigskip
\textbf{b)}\ Berechne die Länge von $\overline{AB}$. Runde auf eine Stelle nach dem Komma.\linie{14mm}\karo{4}}
\fusshilfe{\mbox{5:~a) 5 und 6 Kästchen; b) 7,8}}
\end{pfaufg}

% L3-6: Zusammengesetzt: Pythagoras liefert nur ein Teilstück; das andere Stück kommt dazu (Bild ohne Dreieck)
\begin{pfaufg}{}{6.}{}
Ein Turm besteht aus einem Zylinder und einem Kegeldach. Wie hoch ist der ganze Turm?
\gzwei{\begin{tikzpicture}[line width=0.7pt]\draw (0,2.4) ellipse (0.8 and 0.28);\draw (-0.8,0)--(-0.8,2.4) (0.8,0)--(0.8,2.4);\draw (-0.8,0) arc (180:360:0.8 and 0.28);\draw (-0.8,2.4)--(0,4.0)--(0.8,2.4);\draw[dashed,line width=0.5pt] (0,2.4)--(0.8,2.4);\fill (0,2.4) circle (0.8pt);\node[font=\small,anchor=west] at (0.85,2.4) {$r = 2{,}0$ m};\node[font=\small,anchor=south west] at (0.42,3.25) {$5{,}2$ m};\node[font=\small,anchor=west] at (0.85,1.2) {$18$ m};\end{tikzpicture}}{%
\karo{4}}
\fusshilfe{\mbox{6:~22,8 m}}
\end{pfaufg}

% L3-7: Nur Text: selbst skizzieren, Durchmesser aus dem Radius, Diagonale, Überstand
\begin{pfaufg}{nach P10 ’22}{7.}{}
Ein Glas hat die Form eines Zylinders. Innen ist es $11$ cm hoch und hat den Radius $2{,}8$ cm. Ein Strohhalm steht schräg von der Bodenkante zur gegenüberliegenden oberen Kante. Er ragt $4$ cm heraus. Wie lang ist der Strohhalm? Runde auf eine Stelle nach dem Komma.\par
Skizziere zuerst das Glas mit dem Strohhalm.
\karo{5}
\fusshilfe{\mbox{7:~16,3 cm}}
\end{pfaufg}

\par\vfill\noindent{\footnotesize Vorher: Ist der Winkel recht?\quad\textbullet\quad Weiter: Seiten und Winkel mit sin, cos, tan\par}
\end{document}
